\section{PMD-Mean 的逐 prefix OPD 动力学}
\label{sec:mean-prefix-analysis}

本节固定一个 prefix $x$，在有限词表 $\mathcal V$ 上分析 PMD-Mean 的总体最优解。假设 teacher 与当前 anchor 在每个 token 上均有正概率。每轮 inner 精确求解后执行 hard sync。以下结论描述这一分布空间更新；参数共享、有限 inner、采样误差与 prefix 测度变化需要另外控制。

原始 PMD-Mean 论文的三个相关结论在本地源文件中对应以下标签：
\begin{quote}
总体解定理：\path{thm:pmdmean_solution}。\par
混合正则化解释：\path{eq:mixed_subproblem}。\par
精细损失恒等式：\path{lem:pmdmean_quad_growth}。
\end{quote}
本节采用逐 prefix 的 OPD reward 重新推导这些关系，并进一步证明 teacher KL 下降与 outer 收敛。本文乘子 $\lambda$ 对应未缩放平方损失；若原论文温度为 $\tau=\eta^{-1}$，则 $\lambda=\lambda_{\mathrm{paper}}/\tau^2$。

\subsection{代入 OPD reward 后的 Lambert-$W$ 更新}

固定第 $k$ 轮与 prefix $x$，局部记
\begin{align}
 q_a&=\pi_k(a\mid x), & t_a&=\pi_T(a\mid x),\notag\\
 r_a&=\log\frac{t_a}{q_a}, &
 \kappa&=\KL(q\Vert t)=-\E_q r,\notag\\
 h_a&=r_a+\kappa, & G_a&=\eta h_a.
 \label{eq:mean-local-opd-definitions}
\end{align}
因此 $\E_qh=\E_qG=0$。待求 student 分布记为 $p$，其 log-ratio 为 $U_a=\log(p_a/q_a)$。PMD-Mean 的单轮问题是
\begin{equation}
 \min_{p_a>0,\ \sum_a p_a=1}
 L_{\mathrm{Mean}}(p),
 \qquad
 L_{\mathrm{Mean}}(p)=\frac12\sum_aq_a(U_a-\eta h_a)^2.
 \label{eq:mean-opd-local-problem}
\end{equation}

\begin{theorem}[OPD 下 PMD-Mean 的精确单轮解]
\label{thm:mean-opd-explicit}
式~\eqref{eq:mean-opd-local-problem} 有唯一全局最优分布 $p^\star$。它满足
\begin{align}
 U_a^\star
 &=\eta\left(\log\frac{t_a}{q_a}+\kappa\right)
   -W_0\left(\lambda e^{\eta\kappa}
                      \left(\frac{t_a}{q_a}\right)^\eta\right),\notag\\
 p_a^\star
 &=q_a\exp\left\{\eta\left(\log\frac{t_a}{q_a}+\kappa\right)
   -W_0\left(\lambda e^{\eta\kappa}
                      \left(\frac{t_a}{q_a}\right)^\eta\right)\right\}.
 \label{eq:mean-opd-W-substituted}
\end{align}
这里 $\lambda\ge0$ 由 $\sum_a p_a^\star=1$ 唯一确定，并且
\begin{equation}
 \lambda=\KL(q\Vert p^\star).
 \label{eq:mean-opd-multiplier-kl}
\end{equation}
若 $q\ne t$，则 $\lambda>0$，还可以写成
\begin{equation}
 p_a^\star=\frac{q_a}{\lambda}
 W_0\left(\lambda e^{\eta\kappa}
                \left(\frac{t_a}{q_a}\right)^\eta\right).
 \label{eq:mean-opd-W-ratio}
\end{equation}
若 $q=t$，则 $\lambda=0$ 且 $p^\star=q=t$。
\end{theorem}

\begin{proof}
log-ratio 坐标中的约束是 $\E_qe^U=1$。拉格朗日函数及一阶条件为
\begin{align*}
 \mathscr L(U,\lambda)
 &=\frac12\E_q(U-\eta h)^2+\lambda(\E_qe^U-1),\\
 U_a-\eta h_a+\lambda e^{U_a}&=0.
\end{align*}
令 $v_a=\eta h_a-U_a$，得到 $v_ae^{v_a}=\lambda e^{\eta h_a}$，从而
\[
 U_a=\eta h_a-W_0(\lambda e^{\eta h_a}).
\]
定义归一化函数
\[
 \Phi(\lambda)=\sum_aq_a
 \exp\{\eta h_a-W_0(\lambda e^{\eta h_a})\}.
\]
$W_0$ 在非负实轴上严格递增，因此 $\Phi$ 连续且严格递减。当 $q\ne t$ 时，$h$ 非恒定；严格 Jensen 不等式给出
\[
 \Phi(0)=\E_qe^{\eta h}>e^{\eta\E_qh}=1.
\]
当 $\lambda\to\infty$ 时，恒等式
\[
 e^{\eta h_a-W_0(\lambda e^{\eta h_a})}
 =\lambda^{-1}W_0(\lambda e^{\eta h_a})
\]
说明 $\Phi(\lambda)\to0$。故存在唯一正根。若 $q=t$，则 $h=0$，唯一非负根为零。

对构造出的非负乘子，拉格朗日函数的 Hessian 为
\[
 \nabla_U^2\mathscr L
 =\operatorname{diag}\bigl(q_a(1+\lambda e^{U_a})\bigr)\succ0.
\]
所以一阶条件确定了拉格朗日函数的唯一全局最小点。在任意可行点上，拉格朗日函数等于原目标，故这个可行最小点也是约束问题的唯一全局最优点。

最后对一阶条件取 $q$ 均值，代入 $\E_qh=0$ 与 $\E_qe^{U^\star}=1$：
\[
 \E_qU^\star+\lambda=0,
 \qquad
 \lambda=-\E_q\log\frac{p^\star}{q}=\KL(q\Vert p^\star).
\]
将 $h$ 展开即可得到式~\eqref{eq:mean-opd-W-substituted}。
\end{proof}

恢复全局记号时，$p^\star=\pi_k^{\star,\mathrm{Mean}}(\cdot\mid x)$，$\kappa=\mathcal K_k^T(x)$，而乘子是逐 prefix 的 $\lambda_k(x)$。hard sync 更新为 $\pi_{k+1}=\pi_k^{\star,\mathrm{Mean}}$。式~\eqref{eq:mean-opd-W-substituted} 展示了 teacher--anchor 概率比如何进入 Lambert-$W$；$W_0$ 中的修正随动作变化，不能吸收到一个公共配分常数中。

\subsection{归一化乘子的 R\'enyi 界}

定义 teacher 与 anchor 的混合矩
\begin{equation}
 Z_\alpha(q,t)=\sum_a q_a^{1-\alpha}t_a^\alpha,
 \qquad \alpha>0.
 \label{eq:mean-opd-renyi-Z}
\end{equation}
当 $\alpha\ne1$ 时，R\'enyi 散度满足 $\log Z_\alpha=(\alpha-1)D_\alpha(t\Vert q)$。$Z_1=1$，相关表达式在 $\alpha=1$ 处按连续延拓理解。

\begin{proposition}[OPD 归一化界]
\label{prop:mean-opd-renyi-bounds}
对任意 $\eta>0$，
\begin{equation}
 \frac{Z_\eta(Z_\eta-e^{-\eta\kappa})}{Z_{2\eta}}
 \le\lambda\le\eta\kappa+\log Z_\eta.
 \label{eq:mean-opd-renyi-bounds}
\end{equation}
特别地，$\eta=1$ 时
\begin{equation}
 \frac{1-e^{-\kappa}}{1+\chi^2(t\Vert q)}
 \le\lambda\le\kappa.
 \label{eq:mean-opd-eta-one-bounds}
\end{equation}
\end{proposition}

\begin{proof}
令 $X_a=e^{\eta h_a}$，$A=\E_qX$，$B=\E_qX^2$。式~\eqref{eq:mean-opd-W-ratio} 的归一化条件为
\[
 \lambda=\E_qW_0(\lambda X).
\]
先考虑 $\lambda>0$。利用 $W_0$ 的凹性，
\[
 \lambda\le W_0(\lambda A).
\]
函数 $w\mapsto we^w$ 在非负轴递增，故
\[
 \lambda e^\lambda\le\lambda A,
 \qquad \lambda\le\log A.
\]
另一方面，$W_0(z)\ge z/(1+z)$。此不等式可以由 $z=we^w$ 与 $(1-w)e^w\le1$ 验证；$w\ge1$ 时右侧验证直接成立。再用 Cauchy--Schwarz 得到
\begin{align*}
 \lambda
 &\ge\E_q\frac{\lambda X}{1+\lambda X}\\
 &\ge\frac{(\E_q[\lambda X])^2}
            {\E_q[\lambda X(1+\lambda X)]}
 =\frac{\lambda A^2}{A+\lambda B}.
\end{align*}
消去 $\lambda$ 并整理，得到 $\lambda\ge A(A-1)/B$。当 $\lambda=0$ 时，$q=t$，$A=B=1$，两个界均成立。

OPD reward 的代入给出
\[
 A=e^{\eta\kappa}Z_\eta,
 \qquad B=e^{2\eta\kappa}Z_{2\eta}.
\]
将它们代回即可得到式~\eqref{eq:mean-opd-renyi-bounds}。当 $\eta=1$ 时，$Z_1=1$，且
\[
 Z_2=\sum_a\frac{t_a^2}{q_a}=1+\chi^2(t\Vert q),
\]
因此得到式~\eqref{eq:mean-opd-eta-one-bounds}。
\end{proof}

\subsection{Teacher KL 的单步改进}

\begin{proposition}[精确 Mean 更新的 teacher KL 下降]
\label{prop:mean-opd-kl-descent}
单轮精确解满足
\begin{equation}
 \KL(p^\star\Vert t)
 \le\kappa-\left(\eta^{-1}-1\right)\KL(p^\star\Vert q)
       -\frac{\lambda}{2\eta}\chi^2(p^\star\Vert q).
 \label{eq:mean-opd-kl-descent}
\end{equation}
当 $0<\eta\le1$ 且 $q\ne t$ 时，$\KL(p^\star\Vert t)<\KL(q\Vert t)$。
\end{proposition}

\begin{proof}
固定由 Mean 解确定的 $\lambda$，定义
\[
 J_\lambda(p)=\E_pr-\eta^{-1}\KL(p\Vert q)
                    -\frac{\lambda}{2\eta}\chi^2(p\Vert q).
\]
对 $p_a$ 求导并加入归一化乘子 $\xi$，其一阶条件为
\[
 r_a-\frac{U_a+1}{\eta}
 -\frac{\lambda}{\eta}\left(\frac{p_a}{q_a}-1\right)+\xi=0.
\]
Mean 的 KKT 为 $U_a^\star+\lambda p_a^\star/q_a=\eta(r_a+\kappa)$。取 $\xi=\kappa+(1-\lambda)/\eta$，可验证 $p^\star$ 满足上述条件。$J_\lambda$ 关于 $p$ 的 Hessian 是对角阵，对角项为 $-1/(\eta p_a)-\lambda/(\eta q_a)<0$；因此 $p^\star$ 是其唯一全局最大点。

比较 $p^\star$ 与可行分布 $q$，
\[
 \E_{p^\star}r-\eta^{-1}\KL(p^\star\Vert q)
 -\frac{\lambda}{2\eta}\chi^2(p^\star\Vert q)
 \ge J_\lambda(q)=-\kappa.
\]
再用
\[
 \KL(p^\star\Vert t)=\KL(p^\star\Vert q)-\E_{p^\star}r
\]
得到式~\eqref{eq:mean-opd-kl-descent}。若 $q\ne t$，则 $\lambda>0$。此时 $p^\star\ne q$：否则 KKT 要求 $\eta h_a=\lambda$ 对所有动作成立，而 $h$ 非恒定。这保证 $\chi^2(p^\star\Vert q)>0$，故在 $0<\eta\le1$ 时严格下降。
\end{proof}

\subsection{Log-ratio 跨度与 outer 收敛}

定义逐 prefix 的 log-ratio 跨度
\begin{equation}
 b(q,t)=\max_a\log\frac{q_a}{t_a}
        -\min_a\log\frac{q_a}{t_a}.
 \label{eq:mean-opd-span}
\end{equation}
该量不受公共归一化常数影响。对于正概率分布，$b(q,t)=0$ 当且仅当 $q=t$。

\begin{theorem}[精确 Mean outer 的跨度收敛率]
\label{thm:mean-opd-span-convergence}
若 $0<\eta\le1$，则
\begin{equation}
 b(p^\star,t)
 \le\left(1-\frac{\eta}{1+\eta b(q,t)}\right)b(q,t).
 \label{eq:mean-opd-span-step}
\end{equation}
固定正 teacher、正初始策略和常数 $0<\eta\le1$，每轮精确 Mean inner 后 hard sync。令 $b_k(x)=b(\pi_k(\cdot\mid x),\pi_T(\cdot\mid x))$，则
\begin{equation}
 b_k(x)\le
 \left(1-\frac{\eta}{1+\eta b_0(x)}\right)^k b_0(x)
 \longrightarrow0.
 \label{eq:mean-opd-span-geometric}
\end{equation}
因此 $\pi_k(\cdot\mid x)\to\pi_T(\cdot\mid x)$。当 $\eta=1$ 时，单步还满足
\begin{equation}
 b_{k+1}(x)\le\frac{b_k(x)^2}{1+b_k(x)}.
 \label{eq:mean-opd-span-quadratic}
\end{equation}
\end{theorem}

\begin{proof}
固定本轮 $\lambda,\kappa$，将概率比写成 reward 的连续函数
\[
 w(r)=\exp\{\eta(r+\kappa)-W_0(\lambda e^{\eta(r+\kappa)})\}.
\]
动作 $a$ 的概率比为 $p_a^\star/q_a=w(r_a)$。它满足
\[
 \log w(r)+\lambda w(r)=\eta(r+\kappa),
 \qquad
 \frac{d\log w(r)}{dr}=\frac{\eta}{1+\lambda w(r)}>0.
\]
所以概率比随 reward 严格递增，且其 log-ratio 斜率受到 Lambert-$W$ 修正压缩。对 teacher 的 log-ratio 函数为
\[
 v(r)=\log w(r)-r,
 \qquad v'(r)=-1+\frac{\eta}{1+\lambda w(r)}.
\]
在 $0<\eta\le1$ 时，$v'\le0$。设 $r_{\max}=\max_a r_a$、$r_{\min}=\min_a r_a$，并令 $w_{\max}=w(r_{\max})$。由于 $w$ 递增，
\[
 |v'(r)|\le1-\frac{\eta}{1+\lambda w_{\max}}
 \quad (r_{\min}\le r\le r_{\max}).
\]
均值定理因此给出
\[
 b(p^\star,t)=v(r_{\min})-v(r_{\max})
 \le\left(1-\frac{\eta}{1+\lambda w_{\max}}\right)
       (r_{\max}-r_{\min}).
\]
由归一化 $\E_qw(r)=1$，可知 $w_{\max}\ge1$。其 KKT 与 $\kappa=-\E_qr\le-r_{\min}$ 给出
\[
 \lambda w_{\max}
 =\eta(r_{\max}+\kappa)-\log w_{\max}
 \le\eta(r_{\max}-r_{\min})=\eta b(q,t).
\]
代入即可证明式~\eqref{eq:mean-opd-span-step}。

逐轮应用单步界，先得到 $b_{k+1}\le b_k\le b_0$。函数 $x\mapsto1-\eta/(1+\eta x)$ 递增，因此每轮收缩因子都不超过 $1-\eta/(1+\eta b_0)<1$，归纳得到式~\eqref{eq:mean-opd-span-geometric}。概率归一化同时保证
\[
 \min_a\log\frac{\pi_k(a\mid x)}{\pi_T(a\mid x)}
 \le0\le
 \max_a\log\frac{\pi_k(a\mid x)}{\pi_T(a\mid x)}.
\]
所以每个动作的 log-ratio 绝对值不超过 $b_k(x)$；跨度趋零推出逐动作概率收敛。最后令 $\eta=1$ 即得到式~\eqref{eq:mean-opd-span-quadratic}。
\end{proof}

$\eta=1$ 的二次上界保证超线性收敛；第~\ref{sec:mean-hard-sync-cubic-span} 节进一步利用归一化乘子的二阶性质，得到真实局部三次阶数。一般起点上的一次更新仍是 Lambert-$W$ 分布。不同 prefix 的 $b_0(x)$ 可以不同，因此上述逐 prefix 结论不自动给出全 prefix 的统一收敛率。

\subsection{一般 OPD reward 的大步长极限}

原论文对二元 reward 分析 $\tau\to0$。OPD reward 通常有多个取值，可以直接在其完整取值分布上推广这一极限。以下固定 $q,t$，只令单轮步长 $\eta\to\infty$。

\begin{proposition}[正中心化 log-ratio 的极限分布]
\label{prop:mean-opd-large-eta}
若 $q\ne t$，定义 $[x]_+=\max\{x,0\}$ 与
\[
 c=\sum_aq_a[h_a]_+>0.
\]
则
\begin{equation}
 \frac{\lambda}{\eta}\longrightarrow c,
 \qquad
 p_a^\star\longrightarrow\frac{q_a[h_a]_+}{c}.
 \label{eq:mean-opd-large-eta}
\end{equation}
\end{proposition}

\begin{proof}
$\E_qh=0$ 且 $h$ 非恒定，所以 $h_{\max}>0$，并且 $c>0$。乘子上界给出
\[
 \lambda\le\log\E_qe^{\eta h}\le\eta h_{\max}.
\]
选择满足 $h_a=h_{\max}$ 的动作 $a_+$。概率归一化意味着 $w(r_{a_+})=p_{a_+}^\star/q_{a_+}\le1/q_{a_+}$。于是 KKT 给出
\[
 \eta h_{\max}
 =\log w(r_{a_+})+\lambda w(r_{a_+})
 \le\log(1/q_{a_+})+\lambda/q_{a_+}.
\]
故
\[
 q_{a_+}h_{\max}-\frac{q_{a_+}\log(1/q_{a_+})}{\eta}
 \le\frac{\lambda}{\eta}\le h_{\max}.
\]
$\lambda/\eta$ 的任意序列都有收敛子列，其极限 $c'$ 严格为正。沿该子列，当 $h_a>0$ 时，令 $z_a=\lambda e^{\eta h_a}$。由 $W_0(z)=\log z-\log\log z+o(1)$，
\[
 \frac{W_0(z_a)}{\eta}\to h_a,
 \qquad w(r_a)=\frac{W_0(z_a)}{\lambda}\to\frac{h_a}{c'}.
\]
当 $h_a=0$ 时，$w(r_a)=W_0(\lambda)/\lambda=O(\log\eta/\eta)\to0$。当 $h_a<0$ 时，$z_a\to0$，由 $W_0(z)=z(1+o(1))$ 得到 $w(r_a)=e^{\eta h_a}(1+o(1))\to0$。有限词表允许逐项取极限，归一化因此要求
\[
 1=\lim\sum_aq_aw(r_a)=\frac{\sum_aq_a[h_a]_+}{c'}.
\]
所以所有子列极限都等于 $c$，证明了乘子与分布的整个序列收敛。
\end{proof}

二元 reward $r\in\{0,1\}$、成功率 $p=\E_qr$ 时，$\E_q[r-p]_+=p(1-p)$，因此该一般结论恢复原论文的 $\lambda/\eta\to p(1-p)$。对 OPD，它给出
\[
 p_a^\star\longrightarrow
 \frac{q_a[\log(t_a/q_a)+\KL(q\Vert t)]_+}
 {\E_q[\log(t/q)+\KL(q\Vert t)]_+}.
\]
Gibbs 在大步长下集中到 $t_a/q_a$ 最大的动作，而 Mean 的极限给所有正中心化 reward 的动作分配质量。这个固定 anchor 的大步长极限与 $0<\eta\le1$ 下的 outer teacher 收敛是两个不同结论。

\subsection{有限 inner 的精细损失恒等式}

\begin{proposition}[Mean excess loss 的精确分解]
\label{prop:mean-opd-excess-identity}
对任意正概率分布 $p$，
\begin{equation}
 L_{\mathrm{Mean}}(p)-L_{\mathrm{Mean}}(p^\star)
 =\frac12\E_q\left[\left(\log\frac p{p^\star}\right)^2\right]
       +\lambda\KL(p^\star\Vert p).
 \label{eq:mean-opd-excess-identity}
\end{equation}
\end{proposition}

\begin{proof}
记 $U=\log(p/q)$、$U^\star=\log(p^\star/q)$。展开平方得到
\begin{align*}
 L_{\mathrm{Mean}}(p)-L_{\mathrm{Mean}}(p^\star)
 &=\frac12\E_q(U-U^\star)^2
   +\E_q[(U-U^\star)(U^\star-G)].
\end{align*}
KKT 给出 $U^\star-G=-\lambda p^\star/q$，所以第二项等于
\[
 -\lambda\E_{p^\star}\log\frac p{p^\star}
 =\lambda\KL(p^\star\Vert p).
\]
代回即得所需恒等式。
\end{proof}

例如，总体 inner excess loss 不超过 $\epsilon$ 时，$\E_q[\log(p/p^\star)^2]\le2\epsilon$；若 $q\ne t$，还有 $\KL(p^\star\Vert p)\le\epsilon/\lambda$。接近 teacher 时 $\lambda$ 趋零，后一个界会变松，因此不能据此声称所有轮次都有统一的 KL 误差控制。

\subsection{LOO 中心估计与二元分析的迁移范围}

\begin{proposition}[逐 prefix LOO Mean 的期望损失]
\label{prop:mean-opd-loo-expectation}
固定 prefix、anchor 和候选 student $p$，令 $A_1,\ldots,A_m\overset{\mathrm{iid}}\sim q$，$m\ge2$。定义
\[
 \widehat{\bar r}_{-i}=\frac1{m-1}\sum_{j\ne i}r(A_j),
 \qquad
 \widehat L_{\mathrm{Mean}}(p)
 =\frac1{2m}\sum_{i=1}^m
 \left[U(A_i)-\eta\bigl(r(A_i)-\widehat{\bar r}_{-i}\bigr)\right]^2.
\]
则
\begin{equation}
 \E\widehat L_{\mathrm{Mean}}(p)
 =L_{\mathrm{Mean}}(p)
  +\frac{\eta^2}{2(m-1)}\Var_q[r].
 \label{eq:mean-opd-loo-expectation}
\end{equation}
\end{proposition}

\begin{proof}
令 $\bar r=\E_qr$ 与 $e_{-i}=\widehat{\bar r}_{-i}-\bar r$。第 $i$ 项 residual 为
\[
 U(A_i)-\eta(r(A_i)-\bar r)+\eta e_{-i}.
\]
LOO 保证 $e_{-i}$ 与 $A_i$ 独立，且 $\E e_{-i}=0$、$\E e_{-i}^2=\Var_q[r]/(m-1)$。展开平方取期望，交叉项为零，因此
\[
 \E\left[U(A_i)-\eta(r(A_i)-\widehat{\bar r}_{-i})\right]^2
 =\E_q[U-\eta(r-\bar r)]^2
  +\frac{\eta^2\Var_q[r]}{m-1}.
\]
对 $i$ 求均值并乘 $1/2$ 即得结论。不同 LOO residual 之间的依赖不影响这个期望恒等式。
\end{proof}

附加项与候选 student 无关。因而在有限词表的可微参数化中，冻结 anchor 与 reward 后，LOO Mean 梯度是总体 Mean 梯度的无偏估计。这里的 $\Var_q[r]$ 是完整 prefix OPD log-ratio 的方差。若对整个词表求和，可以直接计算 reward 均值，此时不产生这项中心采样误差。

原论文二元分析中的 $p(1-p)$、成功率 $J(\pi)$ 和 $1-J(\pi)$ 分别依赖 Bernoulli reward 的方差与成功事件解释。OPD 的 reward 通常连续取值且每轮随 anchor 改变，应使用 $\Var_{\pi_k}[\log(\pi_T/\pi_k)]$、teacher KL 与 log-ratio 跨度。总体 Lambert-$W$ 解、混合 KL--$\chi^2$ 解释、归一化界和 excess-loss 恒等式均无需二元假设；二元成功率递推及其统计常数需要重新证明。

式~\eqref{eq:mean-opd-loo-expectation} 是固定候选分布的期望恒等式，不能直接替代数据依赖 ERM 的高概率界。要控制训练得到的 student，仍需给出可实现性或逼近误差、策略类复杂度、优化误差以及 reward/log-ratio 的界。逐 prefix 的有限正概率分布给出有限 reward 范围，但这个范围不一定在所有 prefix 与 outer 轮次上一致。式~\eqref{eq:mean-opd-excess-identity} 则说明：统计分析应围绕约束最优分布的 excess loss 展开；将 Lambert-$W$ 修正本身算作不可消失的 target-estimation 误差，会掩盖 Mean 的总体一致性。
